\section{Cuntz comparison and support calculus}
\label{sec:preliminaries}

We collect the comparison and trace facts used by both routes into
condition (C). Cut placement lets us move positive elements into prescribed
hereditary subalgebras. The trace correspondence then makes numerical
estimates available to the Cuntz semigroup, with zero and infinite values
retained.

All tensor products are spatial unless a maximal tensor product is
explicitly indicated. Unitization is denoted by a tilde. We permit the
adjoined unit in formulas, but maps and constructed elements take
values in the stated algebra.
We write \(\K\) for the compact operators on a separable
infinite-dimensional Hilbert space.

For \(b\in D_+\), set
\[
 \her_D(b)=\overline{bDb},\qquad
 \Ideal_D(b)=\overline{\operatorname{span}}(DbD),\qquad
 p_b=1_{(0,\infty)}(b)\in D^{**}.
\]
The ambient algebra is omitted when clear. Spectral projections are
also used in representations. A comparison or support relation for a
matrix over \(D\) is understood in the corresponding matrix
amplification.

\subsection{Order zero and the Cuntz semigroup}

A completely positive contractive map is abbreviated c.p.c. It has
\emph{order zero} if it preserves orthogonality of positive elements.
For a c.p.c.\ order-zero map \(\theta:E\to D\) with unital domain, the
order-zero structure theorem~\cite{WZ2009} gives
\[
 \theta(z)=h\rho(z)=\rho(z)h,\qquad
 h=\theta(1),\qquad \rho(1)=p_h,
\]
where \(\rho:E\to D^{**}\) is a supporting homomorphism. Its elements
act as multipliers on \(\her(h)\). If \(f:[0,1]\to[0,1]\) is continuous
and \(f(0)=0\), then \(f(h)\rho(z)\) belongs to \(D\) and defines another
c.p.c.\ order-zero map, denoted \(f(\theta)\). A \emph{matrix cone}
means a c.p.c.\ order-zero map from \(M_k\), for some \(k\).

For \(a,b\in(D\otimes\K)_+\), the notation \(a\precsim b\) means that
\(r_j^*br_j\to a\) in norm for suitable \(r_j\in D\otimes\K\).
The Cuntz semigroup \(\Cu(D)\) consists of equivalence classes for the
associated mutual comparison relation, with addition given by
orthogonal sum. We use its usual sequential-supremum structure
\cite{CEI2008}; in particular,
\begin{equation}
\label{eq:pre:cuts}
 [(a-\epsilon)_+]\ll[a],
 \qquad [a]=\sup_{\epsilon>0}[(a-\epsilon)_+].
\end{equation}
Here \(x\ll y\) means that whenever an increasing sequence has
supremum at least \(y\), one of its terms is at least \(x\).
We will use $[a+c]\le[a]+[c]$ and the identities
\begin{equation}
\label{ext:eq:cu-facts}
 [z^*z]=[zz^*],\qquad
 c\in\her(b)_+\Longrightarrow[c]\le[b],\qquad
 [c]=\sup_{n\ge1}[(c-1/n)_+].
\end{equation}
For mutually orthogonal positive elements, the class of the sum is the
sum of the classes: the row of their square roots compares the sum with
its diagonal matrix. Subequivalence to a fixed positive element is closed
under norm limits in the left variable, by choosing a sufficiently good
comparison witness for each approximant.
The semigroup \(W(D\otimes\K)\), defined using finite matrices over the
stable algebra, is \(\Cu(D)\).

We record the elementary comparison facts which will control the
locations and norms of witnesses.

\begin{lemma}[Cuts, placement, and witness bounds]
\label{lem:pre:cut-placement}
Let \(D\) be a \(C^*\)-algebra.
\begin{enumerate}
\item If \(a,b\in D_+\) and \(\|a-r^*br\|<\epsilon\), then there is
\(v\in D\) with
\[
 v^*v=(a-\epsilon)_+,\qquad vv^*\in\her(b).
\]
Consequently, if \(a\precsim b\), the hereditary algebra of every
positive cut of \(a\) can be placed in \(\her(b)\) by a polar
isometry.
\item If \(a\in\Ideal(b)_+\), then for every \(\epsilon>0\) there is
a finite integer \(L\) such that
\[
 [(a-\epsilon)_+]\le L[b].
\]
If \([a]\le L[b]\), every positive cut of \(a\) still compares to
\(L[(b-\delta)_+]\) for a sufficiently small \(\delta>0\).
\item Comparisons between elements of a hereditary subalgebra can
be witnessed in that subalgebra. Comparisons between specified
matrix corners can be witnessed by the appropriate rectangular
matrices.
\item If \(u\precsim(y-\beta)_+\), where \(0\le\alpha<\beta\), then for
each \(\eta>0\) there is \(r\in D\) such that
\[
 \|u-r^*(y-\alpha)_+r\|<\eta,\qquad
 \|r\|^2\le\frac{\|u\|+\eta}{\beta-\alpha}.
\]
\end{enumerate}
\end{lemma}

\begin{proof}
For the first assertion, choose
\(\|a-r^*br\|<\delta<\epsilon\), and a continuous \(g\ge0\), vanishing
on \([0,\delta]\), which is one on \([\epsilon,\|a\|]\).
For \(z=b^{1/2}rg(a)\),
\[
 (\delta-\|a-r^*br\|)\,g(a)^2
 \le z^*z
 \le (\|a\|+\|a-r^*br\|)\,g(a)^2.
\]
Thus \(z^*z\) and \(g(a)^2\) generate the same hereditary algebra,
which contains \((a-\epsilon)_+\). If \(w\) is the polar isometry of
\(z\), conjugation by \(w\) maps \(\her(z^*z)\) into \(\her(b)\).
Moreover \(wx\in D\) for \(x\in\her(z^*z)\), by approximation using
\(w|z|=z\). Taking \(v=w(a-\epsilon)_+^{1/2}\) proves the assertion.
Its matrix version gives the claimed hereditary placement.

For the second assertion, approximate $a^{1/2}$ by a finite sum
$q=\sum_i x_i b y_i$ so closely that $\|a-q^*q\|<\epsilon/2$.
Write $q=XY$, where $X=(x_i b^{1/2})_i$ is a row and
$Y=(b^{1/2}y_i)_i$ is a column. Then
$q^*q\le\|X\|^2Y^*Y$, while $YY^*$ belongs to the hereditary
algebra of $H=\operatorname{diag}(b,\ldots,b)$.
Positive domination and equality of the classes of the two squared
moduli give
\[
 q^*q\precsim Y^*Y\sim YY^*\precsim H.
\]
Choose a comparison witness with error less than $\epsilon/2$ and
apply the first assertion to obtain the claimed cut comparison.
Alternatively, for a fixed comparison witness against
\(\operatorname{diag}(b,\ldots,b)\), replacing \(b\) by
\((b-\delta)_+\) changes the compression by at most
\(\delta\|r\|^2\). This also proves the assertion about an extra target
cut.

Let \(C\subseteq D\) be hereditary and let \((e_i)\) be an approximate
unit of \(C\). Given \(a,b\in C_+\) and a witness \(r\), first choose
an index for which the errors in \(e_i b e_i\) and \(e_i a e_i\) are
small. Then \(e_i r e_i\in C\), and its compression is as close to
\(a\) as desired. This proves the third assertion. The same
compression in matrix algebras gives rectangular witnesses. More
generally, an element whose two squared moduli belong to the
specified hereditary corners lies in their rectangular block, as
is seen by multiplying by approximate units on both sides.

For the last assertion choose \(x\) with
\(\|u-x^*(y-\beta)_+x\|<\eta\), and put
\[
 f(t)=
 \left(\frac{(t-\beta)_+}{(t-\alpha)_+}\right)^{1/2},
 \qquad r=f(y)x,
\]
where the quotient is zero when its denominator vanishes. This is a
continuous function vanishing near zero, and
\[
 f(y)(y-\alpha)_+f(y)=(y-\beta)_+,\qquad
 f(y)^2\le(\beta-\alpha)^{-1}(y-\beta)_+.
\]
The compression and the norm bound follow. If the original
comparison is in the stabilization, compression to the given corner
first gives \(x\in D\).
\end{proof}

\subsection{Extended traces and rank}

We allow all lower semicontinuous extended traces
\(\tau:A_+\to[0,\infty]\), without dense finiteness. Thus \(\tau\) is
additive and positively homogeneous on positives, is lower
semicontinuous in norm, vanishes at zero, and satisfies
\(\tau(x^*x)=\tau(xx^*)\). For \(b\ge0\), write
\[
 d_\tau(b)=\lim_{j\to\infty}\tau(b^{1/j}),
\]
first rescaling to a contraction when necessary. The rank is unchanged
by positive scalar rescaling. Traces
are extended to matrices and to stabilization using the unnormalized
matrix trace: on a positive infinite matrix, the extended trace is the
sum of the traces of its diagonal entries. Additivity and lower
semicontinuity justify these sums, and traciality follows by summing the
squared moduli of the matrix entries.
For a positive contraction \(b\) and \(s>0\),
\begin{equation}
\label{eq:pre:trace-rank}
 \tau(b)\le d_\tau(b),\qquad
 d_\tau((b-s)_+)\le s^{-1}\tau(b).
\end{equation}
These inequalities follow from functional calculus. Rank respects
Cuntz comparison. For a bounded trace we use its normal extension
to the bidual, so \(d_\tau(b)=\tau(p_b)\).

Every such trace defines a functional $[b]\mapsto d_\tau(b)$ on
$\Cu(A)$. Conversely, if $A$ is exact, every
$\varphi\in\F(\Cu(A))$ has the form
\begin{equation}
\label{ext:eq:trace-representation}
 \varphi([b])=d_\tau(b),\qquad b\in(A\otimes\K)_+,
\end{equation}
for a lower semicontinuous extended trace $\tau$ on the stabilization.
Elliott--Robert--Santiago's correspondence first gives an extended
quasitrace; Haagerup's theorem, in its extended form for exact algebras,
makes it additive, and stabilization preserves the trace cone
\cite[Proposition~4.2 and Remarks~4.3 and~4.5]{ERS2011}.
Neither boundedness nor dense finiteness is imposed.

Nuclear algebras are exact. Indeed, for an ideal $J\lhd D$, the quotient
norm on $A\odot(D/J)$ induced by
$(A\otimes_{\min}D)/(A\otimes_{\min}J)$ dominates the spatial norm and
is a $C^*$-cross-norm, as seen by elementary-tensor lifts. Here $\odot$
denotes the algebraic tensor product. Uniqueness of the tensor norm for
nuclear $A$ identifies this quotient with $A\otimes_{\min}(D/J)$,
giving the exact tensor-product sequence.

\subsection{The comparison convention}

The extended-functional convention in
Equation~\eqref{ext:eq:comparison} is equivalent to almost
unperforation. The non-strict inequality on the whole functional cone
records ideal containment; the strict inequalities at finite positive
target values then determine comparison.

\begin{lemma}
\label{lem:comparison-convention}
For any $C^*$-algebra $A$, strict comparison in the sense of
Equation~\eqref{ext:eq:comparison} is equivalent to almost unperforation
of $\Cu(A)$.
\end{lemma}

\begin{proof}
Suppose first that Equation~\eqref{ext:eq:comparison} holds and that
$(l+1)x\le ly$. Applying a functional gives
$\varphi(x)\le\varphi(y)$, with strict inequality whenever
$0<\varphi(y)<\infty$. If $\varphi(y)=0$, then $\varphi(x)=0$;
if $\varphi(y)=\infty$, the non-strict inequality is automatic.
Thus Equation~\eqref{ext:eq:comparison} gives $x\le y$.

Conversely, suppose $\Cu(A)$ is almost unperforated, and let
$x=[a]$ and $y=[b]$ satisfy the two functional inequalities in
Equation~\eqref{ext:eq:comparison}. In $B=A\otimes\K$, set
$I=\Ideal_B(b)$. The map that is zero on $I_+$ and infinite outside
$I_+$ is a lower semicontinuous extended trace: a sum of positives
belongs to an ideal exactly when both summands do, $z^*z$ belongs to it
exactly when $zz^*$ does, and the ideal is norm closed. Its rank
functional vanishes at $y$, so the first inequality forces $a\in I$.
Lemma~\ref{lem:pre:cut-placement}(2) now gives
$x\le\sup_{n\ge1}ny$. For every functional normalized by
$\varphi(y)=1$, the second inequality gives $\varphi(x)<1$.
The comparison criterion of
\cite[Proposition~6.2]{ERS2011} therefore yields $x\le y$.
\end{proof}
